\documentclass[10pt,a4paper]{amsart}
\usepackage{amsmath,amssymb,amsthm,graphicx}
% Retained minimal preamble for the Stacks Project chapter flat.tex (More on Flatness).
% Verbatim excerpt of the chapter preamble: Xy-pic setup (native preamble lines 14-18) and the native
% theorem declarations with the mathematical macros (native preamble lines 154-217). Copied mechanically;
% nothing added, removed or redefined. The wrapper class and the amsmath/amssymb/amsthm/graphicx base
% packages are supplied by the assembler, so the class/externaldocument part of the native preamble is omitted.
% For commutative diagrams we use Xy-pic
\usepackage[all]{xy}

% We use 2cell for 2-commutative diagrams.
\xyoption{2cell}
% Theorem environments.
%
\theoremstyle{plain}
\newtheorem{theorem}[subsection]{Theorem}
\newtheorem{proposition}[subsection]{Proposition}
\newtheorem{lemma}[subsection]{Lemma}

\theoremstyle{definition}
\newtheorem{definition}[subsection]{Definition}
\newtheorem{example}[subsection]{Example}
\newtheorem{exercise}[subsection]{Exercise}
\newtheorem{situation}[subsection]{Situation}

\theoremstyle{remark}
\newtheorem{remark}[subsection]{Remark}
\newtheorem{remarks}[subsection]{Remarks}

\numberwithin{equation}{subsection}

% Macros
%
\def\lim{\mathop{\mathrm{lim}}\nolimits}
\def\colim{\mathop{\mathrm{colim}}\nolimits}
\def\Spec{\mathop{\mathrm{Spec}}}
\def\Hom{\mathop{\mathrm{Hom}}\nolimits}
\def\Ext{\mathop{\mathrm{Ext}}\nolimits}
\def\SheafHom{\mathop{\mathcal{H}\!\mathit{om}}\nolimits}
\def\SheafExt{\mathop{\mathcal{E}\!\mathit{xt}}\nolimits}
\def\Sch{\mathit{Sch}}
\def\Mor{\mathop{\mathrm{Mor}}\nolimits}
\def\Ob{\mathop{\mathrm{Ob}}\nolimits}
\def\Sh{\mathop{\mathit{Sh}}\nolimits}
\def\NL{\mathop{N\!L}\nolimits}
\def\CH{\mathop{\mathrm{CH}}\nolimits}
\def\proetale{{pro\text{-}\acute{e}tale}}
\def\etale{{\acute{e}tale}}
\def\QCoh{\mathit{QCoh}}
\def\Ker{\mathop{\mathrm{Ker}}}
\def\Im{\mathop{\mathrm{Im}}}
\def\Coker{\mathop{\mathrm{Coker}}}
\def\Coim{\mathop{\mathrm{Coim}}}

% Boxtimes
%
\DeclareMathSymbol{\boxtimes}{\mathbin}{AMSa}{"02}

%
% Macros for moduli stacks/spaces
%
\def\QCohstack{\mathcal{QC}\!\mathit{oh}}
\def\Cohstack{\mathcal{C}\!\mathit{oh}}
\def\Spacesstack{\mathcal{S}\!\mathit{paces}}
\def\Quotfunctor{\mathrm{Quot}}
\def\Hilbfunctor{\mathrm{Hilb}}
\def\Curvesstack{\mathcal{C}\!\mathit{urves}}
\def\Polarizedstack{\mathcal{P}\!\mathit{olarized}}
\def\Complexesstack{\mathcal{C}\!\mathit{omplexes}}
% \Pic is the operator that assigns to X its picard group, usage \Pic(X)
% \Picardstack_{X/B} denotes the Picard stack of X over B
% \Picardfunctor_{X/B} denotes the Picard functor of X over B
\def\Pic{\mathop{\mathrm{Pic}}\nolimits}
\def\Picardstack{\mathcal{P}\!\mathit{ic}}
\def\Picardfunctor{\mathrm{Pic}}
\def\Deformationcategory{\mathcal{D}\!\mathit{ef}}

% --- appended 2026-10-12 (Chief): the project own definitions of the reference/slogan/history
% display environments, copied mechanically from the shared Stacks preamble lines 8-12
% (\usepackage{verbatim} plus the three \newenvironment lines, which the chapter preamble keeps
% verbatim from the project preamble). Nothing added, removed or redefined: these environments
% carry no mathematical content and the project itself typesets them as comments. Without these
% three lines a counted statement whose own text uses them fails with Environment slogan undefined.
\usepackage{verbatim}
\newenvironment{reference}{\comment}{\endcomment}
\newenvironment{slogan}{\comment}{\endcomment}
\newenvironment{history}{\comment}{\endcomment}
\usepackage{cleveref}
\crefname{originallemma}{original lemma}{original lemmas}
\Crefname{originallemma}{Original lemma}{Original lemmas}
\makeatletter
\newlabel{algebra-lemma-characterize-valuation-ring}{{090Q}{source}{}{sourcepointer-1}{}}
\newlabel{algebra-lemma-characterize-valuation-ring@cref}{{[originallemma][0][]090Q}{[1][0][]source}{}{}{}}
\makeatother
\crefname{originallemma}{original lemma}{original lemmas}
\Crefname{originallemma}{Original lemma}{Original lemmas}
\makeatletter
\newlabel{algebra-lemma-faithfully-flat-universally-injective}{{05CK}{source}{}{sourcepointer-2}{}}
\newlabel{algebra-lemma-faithfully-flat-universally-injective@cref}{{[originallemma][0][]05CK}{[1][0][]source}{}{}{}}
\makeatother
\crefname{originallemma}{original lemma}{original lemmas}
\Crefname{originallemma}{Original lemma}{Original lemmas}
\makeatletter
\newlabel{algebra-lemma-ideals-valuation-ring}{{00IH}{source}{}{sourcepointer-3}{}}
\newlabel{algebra-lemma-ideals-valuation-ring@cref}{{[originallemma][0][]00IH}{[1][0][]source}{}{}{}}
\makeatother
\crefname{originallemma}{original lemma}{original lemmas}
\Crefname{originallemma}{Original lemma}{Original lemmas}
\makeatletter
\newlabel{algebra-lemma-minimal-contains}{{0AS6}{source}{}{sourcepointer-4}{}}
\newlabel{algebra-lemma-minimal-contains@cref}{{[originallemma][0][]0AS6}{[1][0][]source}{}{}{}}
\makeatother
\crefname{originallemma}{original lemma}{original lemmas}
\Crefname{originallemma}{Original lemma}{Original lemmas}
\makeatletter
\newlabel{lemma-flat-finite-type-local-colimit-free}{{05U7}{source}{}{sourcepointer-5}{}}
\newlabel{lemma-flat-finite-type-local-colimit-free@cref}{{[originallemma][0][]05U7}{[1][0][]source}{}{}{}}
\makeatother
\crefname{originaldefinition}{original definition}{original definitions}
\Crefname{originaldefinition}{Original definition}{Original definitions}
\makeatletter
\newlabel{more-algebra-definition-content-ideal}{{0ASA}{source}{}{sourcepointer-6}{}}
\newlabel{more-algebra-definition-content-ideal@cref}{{[originaldefinition][0][]0ASA}{[1][0][]source}{}{}{}}
\makeatother
\crefname{originallemma}{original lemma}{original lemmas}
\Crefname{originallemma}{Original lemma}{Original lemmas}
\makeatletter
\newlabel{more-algebra-lemma-content-finitely-generated}{{0ASB}{source}{}{sourcepointer-7}{}}
\newlabel{more-algebra-lemma-content-finitely-generated@cref}{{[originallemma][0][]0ASB}{[1][0][]source}{}{}{}}
\makeatother
\crefname{originallemma}{original lemma}{original lemmas}
\Crefname{originallemma}{Original lemma}{Original lemmas}
\makeatletter
\newlabel{more-algebra-lemma-henselization-valuation-ring}{{0ASK}{source}{}{sourcepointer-8}{}}
\newlabel{more-algebra-lemma-henselization-valuation-ring@cref}{{[originallemma][0][]0ASK}{[1][0][]source}{}{}{}}
\makeatother
\title{Content ideals exist and are principal for a flat finite module over an essentially finite type local map out of a valuation ring}
\author{The Stacks Project Authors}
\date{stacks-project:flat.tex:local-snapshot-20261011; GNU FDL 1.2 or later}
\begin{document}
\maketitle
\noindent\textit{Editorial note.} All delivered mathematical text is the author's own, copied from the pinned Stacks Project chapter file flat.tex; no statement or proof step was rewritten, reordered or removed.\par\medskip
\noindent\textit{Editorial note.} Preamble provenance: the chapter's own retained minimal preamble (sources/stacks-flat/retained-preamble.tex, lines 1-85) is a verbatim excerpt of the native Xy-pic setup and of the native theorem declarations and mathematical macros, followed by the project's own reference/slogan/history display-environment definitions copied mechanically from the shared preamble; the wrapper class and the amsmath/amssymb/amsthm/graphicx base packages are supplied by the assembler, and extra\_packages is left empty because everything needed is already in the retained preamble. No macro was invented or redefined.\par\medskip
\noindent\textit{Editorial note.} Same-chapter pointers are located at their \textbackslash{}label lines inside the pinned chapter file; cross-chapter pointers are located in the project's own tag list staged in the snapshot (sources/stacks-chow/tags.txt), at the verbatim lines carrying their canonical tags.\par\medskip
\noindent\textit{Allowed standard prerequisites (editorial).}\begin{itemize}
\item The notion of a content ideal, and its finite generation. More on Algebra, labels more-algebra-definition-content-ideal (statement, proof line 5806) and more-algebra-lemma-content-finitely-generated (proof line 5813); Algebra, label algebra-lemma-characterize-valuation-ring (proof line 5814). Fixes what has to be produced, gives that the content ideal is finitely generated, and with the characterization of a valuation ring shows that a finitely generated content ideal is principal.
\item Faithful flatness of the localisation after henselization and universal injectivity of a faithfully flat map. Algebra, label algebra-lemma-faithfully-flat-universally-injective, invoked at proof line 5828. Shows that the membership of an element in an ideal times the module is detected after the faithfully flat base change, which is the reduction to the henselized situation.
\item Ideals of a valuation ring are in bijection along the henselization, and the content statement over the henselization. More on Algebra, label more-algebra-lemma-henselization-valuation-ring (proof line 5829); Algebra, label algebra-lemma-ideals-valuation-ring (proof line 5830); Same chapter flat.tex, label lemma-flat-finite-type-local-colimit-free (proof line 5835) and Algebra, label algebra-lemma-minimal-contains (proof line 5836). The first two give that ideals correspond bijectively under the henselization so that content ideals transfer, and the last two provide the content statement in the henselized case which is then transported back.
\end{itemize}\medskip
\ifdefined\hypertarget\hypertarget{sourcepointer-1}{}\fi
\noindent\textit{Source reference.} Original lemma 090Q: sources/stacks-chow/tags.txt, source lines 11418--11418. Navigation pointer for \textbackslash{}ref\{algebra-lemma-characterize-valuation-ring\} used inside the counted span (line 5814); the argument that consumes it is retained in the span.\par\medskip
\ifdefined\hypertarget\hypertarget{sourcepointer-2}{}\fi
\noindent\textit{Source reference.} Original lemma 05CK: sources/stacks-chow/tags.txt, source lines 6923--6923. Navigation pointer for \textbackslash{}ref\{algebra-lemma-faithfully-flat-universally-injective\} used inside the counted span (line 5828); the argument that consumes it is retained in the span.\par\medskip
\ifdefined\hypertarget\hypertarget{sourcepointer-3}{}\fi
\noindent\textit{Source reference.} Original lemma 00IH: sources/stacks-chow/tags.txt, source lines 679--679. Navigation pointer for \textbackslash{}ref\{algebra-lemma-ideals-valuation-ring\} used inside the counted span (line 5830); the argument that consumes it is retained in the span.\par\medskip
\ifdefined\hypertarget\hypertarget{sourcepointer-4}{}\fi
\noindent\textit{Source reference.} Original lemma 0AS6: sources/stacks-chow/tags.txt, source lines 13569--13569. Navigation pointer for \textbackslash{}ref\{algebra-lemma-minimal-contains\} used inside the counted span (line 5836); the argument that consumes it is retained in the span.\par\medskip
\ifdefined\hypertarget\hypertarget{sourcepointer-5}{}\fi
\noindent\textit{Source reference.} Original lemma 05U7: sources/stacks-flat/pinned-original-source.tex, source lines 5759--5759. Navigation pointer for \textbackslash{}ref\{lemma-flat-finite-type-local-colimit-free\} used inside the counted span (line 5835); the argument that consumes it is retained in the span.\par\medskip
\ifdefined\hypertarget\hypertarget{sourcepointer-6}{}\fi
\noindent\textit{Source reference.} Original definition 0ASA: sources/stacks-chow/tags.txt, source lines 13573--13573. Navigation pointer for \textbackslash{}ref\{more-algebra-definition-content-ideal\} used inside the counted span (line 5806); the argument that consumes it is retained in the span.\par\medskip
\ifdefined\hypertarget\hypertarget{sourcepointer-7}{}\fi
\noindent\textit{Source reference.} Original lemma 0ASB: sources/stacks-chow/tags.txt, source lines 13574--13574. Navigation pointer for \textbackslash{}ref\{more-algebra-lemma-content-finitely-generated\} used inside the counted span (line 5813); the argument that consumes it is retained in the span.\par\medskip
\ifdefined\hypertarget\hypertarget{sourcepointer-8}{}\fi
\noindent\textit{Source reference.} Original lemma 0ASK: sources/stacks-chow/tags.txt, source lines 13583--13583. Navigation pointer for \textbackslash{}ref\{more-algebra-lemma-henselization-valuation-ring\} used inside the counted span (line 5829); the argument that consumes it is retained in the span.\par\medskip
\begin{lemma}
\label{lemma-flat-finite-type-local-valuation-ring-has-content}
Let $A \to B$ be a local ring map of local rings which is essentially of
finite type. Let $N$ be a finite $B$-module which is flat as an $A$-module.
If $A$ is a valuation ring, then any element of $N$ has a content ideal
$I \subset A$ (More on Algebra, Definition
\ref{more-algebra-definition-content-ideal}). Also, $I$ is a
principal ideal.
\end{lemma}

\begin{proof}
The final statement follows from the fact that $I$ is a finitely
generated ideal by
More on Algebra, Lemma \ref{more-algebra-lemma-content-finitely-generated}
and Algebra, Lemma \ref{algebra-lemma-characterize-valuation-ring}.

\medskip\noindent
Proof of existence of $I$.
Let $A \subset A^h$ be the henselization. Let $B'$ be the localization
of $B \otimes_A A^h$ at the maximal ideal
$\mathfrak m_B \otimes A^h + B \otimes \mathfrak m_{A^h}$.
Then $B \to B'$ is flat, hence faithfully flat.
Let $N' = N \otimes_B B'$.
Let $x \in N$ and let $x' \in N'$ be the image.
We claim that for an ideal $I \subset A$ we have
$x \in IN \Leftrightarrow x' \in IN'$.
Namely, $N/IN \to N'/IN'$ is the tensor product of $B \to B'$
with $N/IN$ and $B \to B'$ is universally injective by
Algebra, Lemma \ref{algebra-lemma-faithfully-flat-universally-injective}.
By More on Algebra, Lemma \ref{more-algebra-lemma-henselization-valuation-ring}
and Algebra, Lemma \ref{algebra-lemma-ideals-valuation-ring}
the map $A \to A^h$ defines an inclusion preserving
bijection $I \mapsto IA^h$ on sets of ideals. We conclude that
$x$ has a content ideal in $A$ if and only if $x'$ has a content
ideal in $A^h$. The assertion for $x' \in N'$ follows from
Lemma \ref{lemma-flat-finite-type-local-colimit-free} and
Algebra, Lemma \ref{algebra-lemma-minimal-contains}.
\end{proof}

\end{document}
